\hypertarget{appendix-9079d3659da4}{}
\section*{De Morgan's laws}
\section{De Morgan's laws}\label{de-morgans-laws}

\textbf{Curriculum depth:} Introduced · \textbf{First seen:} C001 · \textbf{Latest development:} C001

\subsection{Definition and mechanism}\label{definition-and-mechanism}

Negation crossing a conjunction or disjunction swaps the operator: \(\neg(A\land B)=\neg A\lor\neg B\), and dually for OR.

Negating `both' means at least one condition fails; negating `either' means both fail. \texttt{!(visible\ \&\&\ loaded)} equals \texttt{!visible\ \textbar{}\textbar{}\ !loaded}. The identities hold lane by lane for fixed-width vectors. Rewriting comparisons also requires their domain: negating a floating comparison needs care around NaN.

\subsection{Prerequisites}\label{prerequisites}

\begin{itemize}
\tightlist
\item
  \hyperref[appendix-24475e1596a8]{Boolean algebra} --- The logical domain is false/true.
\item
  \hyperref[appendix-8f38927cfc54]{Bitwise NOT} --- Complement must specify width: eight-bit complement of \texttt{0x0F} is \texttt{0xF0}, while sixteen-bit complement is \texttt{0xFFF0}.
\end{itemize}

\subsection{Used by}\label{used-by}

Connect later chunks here when they use this concept.

\subsection{Recall and next depth}\label{recall-and-next-depth}

Explain de morgan's laws without the label, work the example above, then state the assumption whose removal would change the result.

\textbf{Pending:} Apply and reconstruct this concept independently.

\subsection{Source and navigation}\label{source-and-navigation}

\href{Continuous\%20Notes/Computing\%20Foundations\%20-\%20From\%20Bits\%20to\%20C++\%20and\%20Systems\%20APIs\#From\%20a\%20signal\%20to\%20an\%20interpretation}{C001 teaching chapter} · \hyperref[appendix-fd090d982a15]{Part definitions and mechanisms} · \href{}{Dependency Map (02-Foundations-II, appendix-4ad47bdd269a)}
